\documentclass[11pt]{amsart}

\usepackage[T1]{fontenc}
\usepackage{lmodern}
\usepackage{microtype}
\usepackage{mathtools,amssymb,amsthm}
\usepackage[hidelinks]{hyperref}

\hypersetup{
  pdftitle={A Counterexample to a Rainbow-Domination Implication of Zerovnik}
}

\newtheorem{theorem}{Theorem}
\newtheorem{lemma}[theorem]{Lemma}

\title[A rainbow-domination counterexample]
      {A Counterexample to a Rainbow-Domination Implication of \v{Z}erovnik}
\date{}

\subjclass[2020]{05C69}
\keywords{rainbow domination, Tutte--Coxeter graph, cubic graph,
duads and synthemes}

\begin{document}

\begin{abstract}
Let $T$ be the Tutte--Coxeter graph. We prove
\[
  \bigl(\gamma_{r3}(T),\gamma_{r4}(T),\gamma_{r5}(T)\bigr)
  =(16,21,25).
\]
Since $|V(T)|=30$, this gives
\[
  \gamma_{r5}(T)=\frac{5|V(T)|}{6}
  \qquad\text{but}\qquad
  \gamma_{r4}(T)>\frac{2|V(T)|}{3}.
\]
Thus the first implication in Question~2 of \v{Z}erovnik fails in
general for cubic graphs.
\end{abstract}

\maketitle

\section{Introduction}

For a positive integer $k$, write $[k]=\{1,\ldots,k\}$. A
\emph{$k$-rainbow dominating function} ($k$-RDF) on a graph $G$ is a
map $f\colon V(G)\to 2^{[k]}$ such that
\[
  f(v)=\varnothing
  \quad\Longrightarrow\quad
  \bigcup_{u\in N(v)} f(u)=[k].
\]
Its weight is $w(f)=\sum_{v\in V(G)}|f(v)|$, and the minimum weight is
denoted by $\gamma_{rk}(G)$.

Kuzman proved that every cubic graph $G$ of order $N$ satisfies
\begin{equation}\label{eq:bound}
  \gamma_{rk}(G)\geq \left\lceil\frac{kN}{6}\right\rceil,
  \qquad k\leq 6
\end{equation}
\cite[Theorem~1.2]{Kuzman2020}. Following \cite{Kuzman2025}, a cubic
graph $G$ is \emph{$k$-rainbow domination regular} when
$\gamma_{rk}(G)=k|V(G)|/6$, the extreme case of \eqref{eq:bound}.
Writing $N=|V(G)|$, Question~2 of
\v{Z}erovnik asks whether the reverse chain
\[
  \gamma_{r5}(G)=\frac{5N}{6}
  \ \Longrightarrow\
  \gamma_{r4}(G)=\frac{2N}{3}
  \ \Longrightarrow\
  \gamma_{r3}(G)=\frac{N}{2}
\]
holds for every cubic graph $G$, that is, whether a counterexample
exists \cite[Question~2]{Zerovnik2024}.\footnote{In
\cite[Question~2]{Zerovnik2024} the chain is printed as
$\gamma_{r5}(G)=\frac{5}{3}n\Rightarrow\gamma_{r4}(G)=\frac{4}{3}n
\Rightarrow\gamma_{r3}(G)=n$, and $n$ is not redefined at the question;
throughout that paper $n$ is the parameter of the generalized Petersen
graph $P(n,k)$, which is a graph on $2n$ vertices
\cite[p.~443]{Zerovnik2024}. With $N=|V(G)|=2n$ the three printed
equalities are exactly the equality cases $\gamma_{rk}(G)=kN/6$ of
\eqref{eq:bound} for $k=5,4,3$, which is the chain displayed here. This
is also the only non-vacuous reading: assigning $\{1\}$ to every vertex
shows $\gamma_{rk}(G)\leq|V(G)|$ for all $k$ and all $G$, so under the
literal reading $n=|V(G)|$ the hypothesis $\gamma_{r5}(G)=\frac{5}{3}n$
could never be satisfied.}

Kuzman later noted that the Tutte--Coxeter graph $T$ is not
$3$-rainbow domination regular, that is, that
$\gamma_{r3}(T)\neq 3|V(T)|/6=15$ \cite[Question~(Q6),
p.~37]{Kuzman2025}. Since $\gamma_{r5}(T)=5|V(T)|/6$, as the function
$f_5$ of Section~\ref{sec:proof} shows, that remark already refutes the
implication $\gamma_{r5}(G)=\frac{5N}{6}\Rightarrow
\gamma_{r3}(G)=\frac{N}{2}$ obtained by composing the two implications
of Question~2. What is added here is the exact middle value
$\gamma_{r4}(T)=21>20=\frac{2|V(T)|}{3}$, which places the failure
already in the first implication, together with the exact values of all
three parameters.

\begin{theorem}\label{thm:main}
For the Tutte--Coxeter graph $T$,
\[
  \gamma_{r3}(T)=16,\qquad
  \gamma_{r4}(T)=21,\qquad
  \gamma_{r5}(T)=25.
\]
\end{theorem}

Since $|V(T)|=30$, Theorem~\ref{thm:main} refutes the implication
\[
  \gamma_{r5}(G)=\frac{5|V(G)|}{6}
  \quad\Longrightarrow\quad
  \gamma_{r4}(G)=\frac{2|V(G)|}{3}
\]
for cubic graphs $G$ in general: the graph $T$ satisfies the hypothesis
and not the conclusion. The second implication of Question~2 is not
addressed by $T$, which fails its hypothesis; see the Status and scope
section.

\section{The duad--syntheme model}\label{sec:model}

Let $\Omega=[6]$. A \emph{duad} is a two-element subset of $\Omega$,
and a \emph{syntheme} is a partition of $\Omega$ into three duads. Let
$\mathcal D$ and $\mathcal S$ be the sets of duads and synthemes. Both
have size $15$. The graph $T$ is the bipartite incidence graph on
$\mathcal D\sqcup\mathcal S$, with $d\in\mathcal D$ adjacent to
$s\in\mathcal S$ when $d\in s$. Each duad lies in three synthemes, so
$T$ is cubic of order $30$. The incidence structure is self-dual, and
$T$ is connected; indeed $T$ is the Tutte $8$-cage, the unique cubic
graph of girth $8$ on $30$ vertices \cite{Coxeter1950}. Being connected
and bipartite, $T$ has exactly one bipartition, namely
$\{\mathcal D,\mathcal S\}$.

For either bipartition class $X$, call $Q\subseteq X$ a
\emph{perfect class} if every vertex in the opposite class has exactly
one neighbor in $Q$.

\begin{lemma}\label{lem:cover}
At least five perfect classes are required to cover either bipartition
class of $T$.
\end{lemma}

\begin{proof}
First take $Q\subseteq\mathcal D$. Counting incidences with
$\mathcal S$ gives $3|Q|=15$, so $|Q|=5$. No two members of $Q$ are
disjoint, since two disjoint duads belong to a common syntheme. Hence
$Q$ is a pairwise-intersecting family of five edges of $K_6$, and
therefore a full star. Indeed, if such a family had no common endpoint,
three of its edges would form a triangle, which admits no fourth edge
meeting all three.

Conversely, for each $i\in\Omega$, the star
\[
  \mathcal D_i=\{d\in\mathcal D:i\in d\}
\]
is a perfect class. Thus the perfect classes on $\mathcal D$ are
precisely the six stars $\mathcal D_i$. Four stars leave uncovered the
duad joining two unchosen centers, whereas five stars cover
$\mathcal D$. Self-duality gives the same conclusion on $\mathcal S$.
\end{proof}

\section{Proof of the theorem}\label{sec:proof}

We first record the equality structure needed below. Its first assertion
is the case $d=3$, $n=30$ of \cite[Theorem~3.1]{Kuzman2020}, restated as
\cite[Theorem~2.1(ii)]{Kuzman2025}; the short proof is included because
we also need the assertion about the individual color supports.

\begin{lemma}\label{lem:equality}
Let $G$ be a cubic graph of order $30$, let $k<6$, and let $f$ be a
$k$-RDF of weight $5k$. Then the set of empty and the set of nonempty
vertices are independent sets of size $15$, so they are the two classes
of a bipartition of $G$. Moreover, each color occurs on exactly
five vertices, and every empty vertex has exactly one neighbor carrying
each color.
\end{lemma}

\begin{proof}
Let
\[
  C=\{v\in V(G):f(v)\neq\varnothing\},\qquad c=|C|,
\]
and let $e_C$ be the number of edges in $G[C]$. Count triples
$(x,y,i)$ such that $x\notin C$, $xy\in E(G)$, and $i\in f(y)$. If
$\tau$ is their number, then rainbow domination gives
\[
  \tau\geq k(30-c).
\]
On the other hand,
\[
\begin{aligned}
  \tau
  &=\sum_{y\in C}|f(y)|\bigl(3-\deg_C(y)\bigr)\\
  &\leq 3w(f)-2e_C
   =15k-2e_C.
\end{aligned}
\]
It follows first that $c\geq15$. The cut from $C$ to its complement
has size $3c-2e_C$ and at most $3(30-c)$, so
\[
  e_C\geq3(c-15).
\]
Consequently,
\[
  k(30-c)\leq15k-6(c-15),
\]
or equivalently $(6-k)(c-15)\leq0$. Since $k<6$, we obtain $c=15$.
The preceding inequalities then force $e_C=0$. All $45$ edges incident
with $C$ therefore join $C$ to its complement, so the complement is
also independent.

For a color $i$, let $Q_i=\{v:i\in f(v)\}$. Every vertex outside $C$
has a neighbor in $Q_i$, hence $3|Q_i|\geq15$. Since
\[
  \sum_{i=1}^k |Q_i|=w(f)=5k,
\]
we have $|Q_i|=5$ for every $i$. Its $15$ incident edges meet all $15$
vertices outside $C$, so each such vertex has exactly one neighbor in
$Q_i$.
\end{proof}

\begin{proof}[Proof of Theorem~\ref{thm:main}]
In each construction below, every syntheme receives the empty label and
every duad receives a nonempty label. It is therefore enough to check
the syntheme vertices. Nonemptiness of the duad labels is immediate in
each of the three cases; note that $\{5,6\}$ is the only duad disjoint
from $[4]$, which is exactly why the definition of $f_4$ needs one
exceptional value.

For five colors, set
\[
  f_5(d)=d\cap[5],\qquad d\in\mathcal D.
\]
Because each syntheme partitions $[6]$, the union of its three labels
is $[5]$. Each color occurs on five duads, so $w(f_5)=25$. Together
with \eqref{eq:bound}, this gives $\gamma_{r5}(T)=25$.

For four colors, set
\[
  f_4(d)=
  \begin{cases}
    \{1\},&d=\{5,6\},\\
    d\cap[4],&d\neq\{5,6\}.
  \end{cases}
\]
For every syntheme $s$,
\[
  \bigcup_{d\in s}(d\cap[4])=[4],
\]
and the exceptional label only adds a color. Thus $f_4$ is a $4$-RDF
of weight $4\cdot5+1=21$.

For three colors, set
\[
  f_3(d)=
  \begin{cases}
    \{1,2\},&d=\{1,6\},\\
    \{1\},&6\in d\text{ and }1\notin d,\\
    \{2\},&1\in d\text{ and }6\notin d,\\
    \{3\},&d\cap\{1,6\}=\varnothing.
  \end{cases}
\]
If a syntheme contains $\{1,6\}$, that duad supplies colors $1,2$ and
the other two supply color $3$. Otherwise, the duads containing $6$
and $1$ supply colors $1$ and $2$, while the third supplies color $3$.
Hence $f_3$ is a $3$-RDF. Its three color supports have sizes $5$, $5$,
and $\binom42=6$, so $w(f_3)=16$.

It remains to exclude weights $15$ and $20$. Suppose that a $k$-RDF
of weight $5k$ exists for $k\in\{3,4\}$. By Lemma~\ref{lem:equality},
the empty and the nonempty vertices are the two classes of a bipartition
of $T$; as $T$ is connected, that bipartition is the unique one, so the
nonempty vertices form $\mathcal D$ or $\mathcal S$. Each color support
is therefore a perfect class in that bipartition class, and the $k$
supports cover it. This contradicts
Lemma~\ref{lem:cover}, since $k<5$. Therefore
\[
  \gamma_{r3}(T)\geq16,
  \qquad
  \gamma_{r4}(T)\geq21.
\]
The displayed constructions attain these bounds.
\end{proof}

\section*{Status and scope}

Theorem~\ref{thm:main} refutes the first implication of Question~2 of
\cite{Zerovnik2024} in general for cubic graphs: $T$ is cubic,
$\gamma_{r5}(T)=25=5|V(T)|/6$, and $\gamma_{r4}(T)=21>20=2|V(T)|/3$. It
does not refute the second implication
$\gamma_{r4}(G)=\frac{2|V(G)|}{3}\Rightarrow
\gamma_{r3}(G)=\frac{|V(G)|}{2}$, whose hypothesis $T$ fails, and we make
no claim about that implication. Nothing here characterizes the cubic
graphs for which \eqref{eq:bound} is tight for a given $k$, and no claim
is made about graphs other than $T$. The chain refuted is the one printed
in \cite[Question~2]{Zerovnik2024} after the substitution $N=2n$
recorded in the footnote of Section~1; that substitution is forced,
because the printed equalities are the equality cases of
\eqref{eq:bound}.

The composed implication was already refutable from the literature.
Kuzman's remark that $T$ is not $3$-rainbow domination regular
\cite[Question~(Q6), p.~37]{Kuzman2025}, together with
$\gamma_{r5}(T)=5|V(T)|/6$, refutes
$\gamma_{r5}(G)=\frac{5N}{6}\Rightarrow\gamma_{r3}(G)=\frac{N}{2}$, so a
counterexample to the chain read end to end was available before this
note. What is added here is the exact middle value $\gamma_{r4}(T)=21$,
which locates the failure already in the first implication, together
with the exact values $\gamma_{r3}(T)=16$ and $\gamma_{r5}(T)=25$. The
first assertion of Lemma~\ref{lem:equality} is likewise not new: it is
the case $d=3$, $n=30$ of \cite[Theorem~3.1]{Kuzman2020}, restated as
\cite[Theorem~2.1(ii)]{Kuzman2025}, and is reproved here only to keep
the argument self-contained and to record the statement about individual
color supports.

\medskip
\noindent\textbf{Exact verification.}
Each of the finite checks enumerated below is recomputable from the data
printed in this paper, and each is hand-sized. The only input is the
duad--syntheme model of Section~\ref{sec:model}, which specifies the $15$
duads and $15$ synthemes of $\Omega=[6]$ completely, together with the
three labelings
$f_5$, $f_4$, $f_3$ of Section~\ref{sec:proof}. From these one writes out
the fifteen synthemes, evaluates each labeling
on the three duads of each syntheme and confirms that the union of the
three labels is $[5]$, $[4]$ and $[3]$ respectively, that every duad
label is nonempty, and that the weights are $25$, $21$ and $16$;
Lemma~\ref{lem:cover} then needs only the fact that the
pairwise-intersecting five-member families of duads are the six stars,
and Lemma~\ref{lem:equality} is a counting argument with no computation
in it. The checks were re-done independently, from the printed
definitions rather than from the first pass, and all $93$ of them --- the
$45$ syntheme unions, the $45$ nonemptiness conditions and the three
weight sums --- agreed with the values stated above. The identification of
$T$ as the Tutte $8$-cage is of a different kind: it rests on the
uniqueness of the cubic graph of girth $8$ on $30$ vertices, which is
quoted from \cite{Coxeter1950}, is not hand-sized, and is not among those
$93$ checks; it serves only to name $T$, and no step of the proof uses the
girth.

At this size no program is required, and the $93$ checks just listed were
carried out by hand. They were nevertheless also run by machine, as a
cross-check on that hand work. The incidence structure was rebuilt from
the definitions of Section~\ref{sec:model} and found to be cubic of order
$30$, connected and of girth $8$; the $45$ syntheme unions and the $45$
nonemptiness conditions were re-evaluated, with no failing syntheme and no
empty duad label for any of $f_5$, $f_4$, $f_3$; the perfect classes on
the duad side were enumerated exhaustively, returning exactly the six
stars $\mathcal D_i$ and least cover size five, as in
Lemma~\ref{lem:cover}; and the minimum of $\sum_{v,i}x_{v,i}$ over
$x\in\{0,1\}^{V(T)\times[k]}$ subject to
\[
  \sum_{u\in N(v)}x_{u,i}+\sum_{j\in[k]}x_{v,j}\geq1
  \qquad (v\in V(T),\ i\in[k]),
\]
an integer program whose feasible points are exactly the $k$-RDFs of $T$
under $f(v)=\{i:x_{v,i}=1\}$ and whose optimum is therefore
$\gamma_{rk}(T)$, was computed to optimality as $16$, $21$ and $25$ for
$k=3,4,5$. A verification program accompanies this note and repeats the
cross-checks just described, with one difference: it invokes no solver, and
obtains $16$, $21$ and $25$ from the displayed labelings and from the lower
bounds recalled in the next sentence. No part of the argument rests on
them: the upper bounds are the three displayed labelings, and the lower
bounds are \eqref{eq:bound} together with Lemmas~\ref{lem:cover}
and~\ref{lem:equality}.

A search of the rainbow-domination literature, including
\cite{Kuzman2020,Kuzman2025,Zerovnik2024}, turned up the question and
Kuzman's remark, but no determination of $\gamma_{r4}(T)$ and no
counterexample to the first implication. That search was not exhaustive
and has not been through external review, so any earlier resolution we
have missed takes precedence over this note.

\begin{thebibliography}{9}

\bibitem{Coxeter1950}
H.~S.~M. Coxeter,
\emph{Self-dual configurations and regular graphs},
Bull. Amer. Math. Soc. \textbf{56} (1950), 413--455.

\bibitem{Kuzman2020}
B. Kuzman,
\emph{On $k$-rainbow domination in regular graphs},
Discrete Appl. Math. \textbf{284} (2020), 454--464,
\href{https://doi.org/10.1016/j.dam.2020.04.003}
{doi:10.1016/j.dam.2020.04.003}.

\bibitem{Kuzman2025}
B. Kuzman,
\emph{On cubic rainbow domination regular graphs},
Discrete Appl. Math. \textbf{373} (2025), 26--38,
\href{https://doi.org/10.1016/j.dam.2025.04.046}
{doi:10.1016/j.dam.2025.04.046}.

\bibitem{Zerovnik2024}
J. \v{Z}erovnik,
\emph{On rainbow domination of generalized Petersen graphs $P(ck,k)$},
Discrete Appl. Math. \textbf{357} (2024), 440--448,
\href{https://doi.org/10.1016/j.dam.2024.07.052}
{doi:10.1016/j.dam.2024.07.052}.

\end{thebibliography}

\end{document}